\section{Descripción de protocolos y aplicaciones}

La aplicación servidor actúa como intermediario entre clientes SOCKS5 y
servicios TCP. En el mismo proceso ofrece un plano de administración separado,
al que accede la aplicación cliente de terminal. Esta división evita que las
operaciones administrativas se confundan con el tráfico que retransmite el
proxy.

\subsection{SOCKS5}

El proxy acepta autenticación usuario/contraseña y el comando \texttt{CONNECT}
para destinos IPv4, IPv6 o FQDN. Después de establecer la conexión saliente,
retransmite bytes en ambos sentidos sin interpretar el protocolo de aplicación.
Si un nombre resuelve a varias direcciones, conserva la lista y prueba la
siguiente cuando una conexión falla. La decisión importante no fue reproducir
los pasos de RFC1928 y RFC1929, sino integrarlos con lecturas parciales y con una
máquina de estados que nunca presupone cómo se segmentan los datos en TCP.

\subsection{Protocolo de monitoreo}

El Protocolo de Monitoreo y Configuración (PMC) usa una conexión TCP persistente.
Es textual y orientado a líneas terminadas en \texttt{CRLF}. Cada pedido empieza
con un nombre de comando hasta el primer espacio o fin de línea; los campos
restantes se separan por un único espacio. Para cada pedido el servidor devuelve
una respuesta que comienza con \texttt{+OK} o \texttt{-ERR}. Antes de operar, el
cliente negocia la versión y autentica una credencial administrativa.

Los comandos implementados son \texttt{ADD-USER}, \texttt{DEL-USER},
\texttt{LIST-USERS}, \texttt{METRICS}, \texttt{GET-CONFIG},
\texttt{SET-CONFIG} y \texttt{QUIT}. Las respuestas multi-línea usan prefijo de
cantidad para delimitar la respuesta sin depender del cierre de la conexión. La
gramática, los estados y cada respuesta posible se especifican en
\texttt{docs/mgmt-protocol-rfc.md}.

El cliente de terminal oculta la negociación y el formato de cable. Traduce
subcomandos como \texttt{metrics} o \texttt{add-user} a pedidos PMC y presenta
resultados etiquetados, listas y mensajes de error, en lugar de copiar las líneas
recibidas del socket.
